\section{2.2
主要敏捷方法论}\label{ux4e3bux8981ux654fux6377ux65b9ux6cd5ux8bba}

\begin{quote}
本小节是第二节''敏捷开发入门''的一部分
\end{quote}

\subsection{导航}\label{ux5bfcux822a}

\begin{itemize}
\tightlist
\item
  上一小节: \href{section03-02-01.md}{2.1 敏捷开发概述}
\item
  下一小节: \href{section03-02-03.md}{2.3
  敏捷实践在智慧水利项目中的应用}
\end{itemize}

在了解了敏捷开发的基本概念和原则后，本小节将详细介绍三种主流的敏捷方法论：Scrum、看板方法（Kanban）和极限编程（XP）。这些方法各有侧重点，但都遵循敏捷宣言的核心价值观，可以根据智慧水利平台的项目特点灵活选择和组合使用。

\subsection{2.2.1 Scrum框架}\label{scrumux6846ux67b6}

Scrum是目前应用最广泛的敏捷方法论之一，它提供了一个简单而有效的框架，帮助团队以迭代和增量的方式交付复杂产品。Scrum特别适合需求变化较频繁、技术复杂度高的智慧水利平台项目。

\subsubsection{Scrum的核心角色}\label{scrumux7684ux6838ux5fc3ux89d2ux8272}

Scrum定义了三个核心角色，各自承担不同责任：

\begin{enumerate}
\def\labelenumi{\arabic{enumi}.}
\tightlist
\item
  \textbf{产品负责人（Product Owner）}：

  \begin{itemize}
  \tightlist
  \item
    代表业务方和用户利益，明确产品方向和优先级
  \item
    定义和维护产品待办事项列表（Product Backlog）
  \item
    确保团队理解需求并按价值优先开发
  \item
    在智慧水利平台中，通常由资深水利专家或业务负责人担任，如水文分析负责人、水资源管理专家等
  \end{itemize}
\item
  \textbf{开发团队（Development Team）}：

  \begin{itemize}
  \tightlist
  \item
    跨职能自组织团队，通常由5-9人组成
  \item
    集体负责实现产品功能和质量保证
  \item
    具有自主决定如何完成工作的权力
  \item
    在智慧水利平台中，可能包括前端开发、后端开发、数据库专家、GIS工程师、水文模型专家等不同角色
  \end{itemize}
\item
  \textbf{Scrum主管（Scrum Master）}：

  \begin{itemize}
  \tightlist
  \item
    确保团队理解并遵循Scrum规则和实践
  \item
    帮助团队移除障碍，促进高效工作
  \item
    协助产品负责人有效管理产品待办事项
  \item
    促进团队内部及与外部相关方的沟通协作
  \item
    在智慧水利平台中，通常由了解水利领域且熟悉敏捷方法的人员担任
  \end{itemize}
\end{enumerate}

\subsubsection{Scrum的工件}\label{scrumux7684ux5de5ux4ef6}

Scrum定义了三个主要工件，用于提高透明度和减少沟通成本：

\begin{enumerate}
\def\labelenumi{\arabic{enumi}.}
\tightlist
\item
  \textbf{产品待办事项列表（Product Backlog）}：

  \begin{itemize}
  \tightlist
  \item
    包含所有待开发功能、修复缺陷和技术改进的有序列表
  \item
    由产品负责人维护，不断演进和调整
  \item
    条目按业务价值、风险、优先级和依赖关系排序
  \item
    在智慧水利平台中，可能包括''实现水库调度方案自动生成功能''、``集成遥感数据分析模块''等条目
  \end{itemize}
\item
  \textbf{冲刺待办事项列表（Sprint Backlog）}：

  \begin{itemize}
  \tightlist
  \item
    当前迭代（冲刺）计划完成的工作列表
  \item
    包含从产品待办事项中选择的高优先级条目和实现计划
  \item
    由开发团队维护，动态反映剩余工作情况
  \item
    在智慧水利平台中，一个冲刺待办事项可能包括''完成水位数据可视化界面''和相关任务
  \end{itemize}
\item
  \textbf{产品增量（Increment）}：

  \begin{itemize}
  \tightlist
  \item
    冲刺结束时交付的产品功能集合
  \item
    必须达到''完成''状态，可用于演示或部署
  \item
    与之前的增量集成，形成完整的产品
  \item
    在智慧水利平台中，可能是一个可以演示的新功能模块，如''洪水预警分析子系统V1.0''
  \end{itemize}
\end{enumerate}

\subsubsection{Scrum的事件}\label{scrumux7684ux4e8bux4ef6}

Scrum定义了五个正式事件，用于创建规律性并减少不必要的会议：

\begin{enumerate}
\def\labelenumi{\arabic{enumi}.}
\tightlist
\item
  \textbf{冲刺（Sprint）}：

  \begin{itemize}
  \tightlist
  \item
    固定时长的开发周期，通常为2-4周
  \item
    每个冲刺都有明确的目标和交付物
  \item
    在冲刺期间，不应改变目标或降低质量标准
  \item
    在智慧水利平台中，可以根据项目复杂度和团队情况选择适当的冲刺长度
  \end{itemize}
\item
  \textbf{冲刺规划会议（Sprint Planning）}：

  \begin{itemize}
  \tightlist
  \item
    冲刺开始时举行，确定本次冲刺要完成的工作
  \item
    产品负责人说明高优先级条目，团队决定能够承诺完成的范围
  \item
    输出冲刺目标和冲刺待办事项列表
  \item
    在智慧水利平台中，可能讨论''本次冲刺我们将完成水文监测数据实时采集和基础展示功能''
  \end{itemize}
\item
  \textbf{每日站会（Daily Scrum）}：

  \begin{itemize}
  \tightlist
  \item
    每天固定时间举行的简短会议（通常15分钟）
  \item
    团队成员分享昨天完成的工作、今天计划和遇到的障碍
  \item
    关注进度和协作，不是详细的问题解决讨论
  \item
    在智慧水利平台团队中，可能需要特别关注跨学科协作问题
  \end{itemize}
\item
  \textbf{冲刺评审会（Sprint Review）}：

  \begin{itemize}
  \tightlist
  \item
    冲刺结束时举行，向相关方展示完成的功能
  \item
    获取反馈并调整产品待办事项
  \item
    关注产品功能和业务价值
  \item
    在智慧水利平台中，可邀请水利部门用户参与，实际操作和评价新功能
  \end{itemize}
\item
  \textbf{冲刺回顾会（Sprint Retrospective）}：

  \begin{itemize}
  \tightlist
  \item
    冲刺评审后举行，团队反思过程和协作
  \item
    识别做得好的方面和需要改进的地方
  \item
    制定具体的改进计划
  \item
    在智慧水利平台团队中，可能讨论如何改进水利专家与技术人员的协作方式
  \end{itemize}
\end{enumerate}

\subsubsection{Scrum在智慧水利平台中的应用}\label{scrumux5728ux667aux6167ux6c34ux5229ux5e73ux53f0ux4e2dux7684ux5e94ux7528}

Scrum方法在智慧水利平台开发中具有以下优势：

\begin{enumerate}
\def\labelenumi{\arabic{enumi}.}
\item
  \textbf{灵活应对变化}：水利信息化需求常受政策调整、极端气候事件等因素影响，Scrum的迭代式开发能够及时响应这些变化。
\item
  \textbf{促进跨学科协作}：智慧水利平台需要水利、软件、数据科学等多领域专家协作，Scrum的团队结构和日常实践有助于打破专业壁垒。
\item
  \textbf{降低技术风险}：水利平台常涉及复杂算法和模型，Scrum的短周期验证可以及早发现技术问题。
\item
  \textbf{提高用户满意度}：通过频繁展示和获取反馈，确保系统真正满足水利行业用户的实际需求。
\end{enumerate}

案例：某省防汛抗旱指挥系统采用Scrum方法开发，通过3周一个冲刺的节奏，先后完成了基础数据管理、实时监测、预警预报、辅助决策等核心功能。在2021年汛期到来前，系统已经可以投入试运行，及时支持了当年的防汛工作。与传统方法相比，采用Scrum开发的系统更符合一线防汛人员的使用习惯，功能更加贴合实际需求。

\subsection{2.2.2
看板方法（Kanban）}\label{ux770bux677fux65b9ux6cd5kanban}

看板方法源于丰田生产系统，是一种可视化工作流程管理方法，强调持续交付而非固定周期迭代。看板特别适合运维类项目和需求流较平稳的智慧水利平台子系统开发。

\subsubsection{看板的核心原则}\label{ux770bux677fux7684ux6838ux5fc3ux539fux5219}

看板方法基于以下五个核心原则：

\begin{enumerate}
\def\labelenumi{\arabic{enumi}.}
\tightlist
\item
  \textbf{可视化工作流程}：

  \begin{itemize}
  \tightlist
  \item
    使用看板墙展示所有工作项及其状态
  \item
    清晰展示工作流中的瓶颈和阻塞
  \item
    在智慧水利平台中，可以针对不同类型工作（如新功能开发、缺陷修复、模型优化等）使用不同颜色的卡片
  \end{itemize}
\item
  \textbf{限制在制品数量（WIP）}：

  \begin{itemize}
  \tightlist
  \item
    为每个工作阶段设置数量限制
  \item
    确保团队专注完成已开始的工作，避免多任务切换
  \item
    在智慧水利平台中，可能需要根据团队技能分布和任务复杂度调整不同阶段的WIP限制
  \end{itemize}
\item
  \textbf{管理工作流程}：

  \begin{itemize}
  \tightlist
  \item
    监控和优化工作流程，消除瓶颈
  \item
    关注周期时间（从开始到完成的时间）
  \item
    在智慧水利平台中，特别关注跨学科协作环节的流转效率
  \end{itemize}
\item
  \textbf{明确流程规则}：

  \begin{itemize}
  \tightlist
  \item
    团队共同定义清晰的工作规则
  \item
    包括任务分类、优先级确定、质量标准等
  \item
    在智慧水利平台中，需要明确设定水利业务验收标准
  \end{itemize}
\item
  \textbf{持续改进}：

  \begin{itemize}
  \tightlist
  \item
    通过数据分析和定期回顾不断优化流程
  \item
    鼓励团队实验和创新
  \item
    在智慧水利平台中，可以利用周期时间和吞吐量数据优化开发流程
  \end{itemize}
\end{enumerate}

\subsubsection{看板的实施步骤}\label{ux770bux677fux7684ux5b9eux65bdux6b65ux9aa4}

\begin{enumerate}
\def\labelenumi{\arabic{enumi}.}
\tightlist
\item
  \textbf{可视化当前流程}：

  \begin{itemize}
  \tightlist
  \item
    识别工作类型和工作流阶段
  \item
    创建看板墙，反映实际工作状态
  \item
    在智慧水利平台中，典型的流程可能包括：需求分析→设计→开发→测试→部署→验收
  \end{itemize}
\item
  \textbf{应用WIP限制}：

  \begin{itemize}
  \tightlist
  \item
    分析团队能力和工作特点
  \item
    为各阶段设置合理的WIP限制
  \item
    根据实际情况调整限制
  \item
    在智慧水利平台中，开发阶段的WIP可能需要根据技术复杂度设置
  \end{itemize}
\item
  \textbf{管理工作流}：

  \begin{itemize}
  \tightlist
  \item
    每日站会讨论工作流动情况
  \item
    主动解决阻塞问题
  \item
    优先完成已开始的工作
  \item
    在智慧水利平台中，特别关注依赖外部数据源或系统的任务进展
  \end{itemize}
\item
  \textbf{收集并分析指标}：

  \begin{itemize}
  \tightlist
  \item
    周期时间：完成一项工作的总时间
  \item
    吞吐量：单位时间内完成的工作项数量
  \item
    质量指标：缺陷率、客户满意度等
  \item
    在智慧水利平台中，可以分析不同类型功能（如数据处理、模型计算、可视化等）的周期时间差异
  \end{itemize}
\item
  \textbf{持续优化}：

  \begin{itemize}
  \tightlist
  \item
    定期回顾会议分析流程数据
  \item
    识别改进机会并实施改进措施
  \item
    评估改进效果
  \item
    在智慧水利平台中，可能需要根据季节性业务变化（如汛期/非汛期）调整工作流程
  \end{itemize}
\end{enumerate}

\subsubsection{看板与Scrum的区别和结合使用}\label{ux770bux677fux4e0escrumux7684ux533aux522bux548cux7ed3ux5408ux4f7fux7528}

\begin{longtable}[]{@{}
  >{\raggedright\arraybackslash}p{(\columnwidth - 6\tabcolsep) * \real{0.1579}}
  >{\raggedright\arraybackslash}p{(\columnwidth - 6\tabcolsep) * \real{0.1579}}
  >{\raggedright\arraybackslash}p{(\columnwidth - 6\tabcolsep) * \real{0.1842}}
  >{\raggedright\arraybackslash}p{(\columnwidth - 6\tabcolsep) * \real{0.5000}}@{}}
\toprule\noalign{}
\begin{minipage}[b]{\linewidth}\raggedright
特性
\end{minipage} & \begin{minipage}[b]{\linewidth}\raggedright
看板
\end{minipage} & \begin{minipage}[b]{\linewidth}\raggedright
Scrum
\end{minipage} & \begin{minipage}[b]{\linewidth}\raggedright
智慧水利平台应用场景
\end{minipage} \\
\midrule\noalign{}
\endhead
\bottomrule\noalign{}
\endlastfoot
迭代周期 & 持续流动，无固定周期 & 固定时长的冲刺 &
看板：水文数据处理等持续性运维工作Scrum：防汛决策系统等功能开发 \\
角色定义 & 无严格角色划分 & 明确的三个角色 &
看板：成熟团队的自组织管理Scrum：需要明确责任划分的新团队 \\
变更管理 & 随时可变更 & 冲刺内保持稳定 &
看板：频繁变化的监测站点管理Scrum：相对稳定的预报模型开发 \\
工作量控制 & 限制WIP & 基于团队速度的冲刺规划 &
看板：根据团队容量限制并行任务数Scrum：根据历史数据规划迭代工作量 \\
会议要求 & 灵活，通常只有每日站会 & 五个正式事件 &
看板：分布式团队的日常协作Scrum：需要严格过程控制的关键系统开发 \\
\end{longtable}

在智慧水利平台开发中，可以结合使用这两种方法，如采用''Scrumban''：在开发新功能时使用Scrum的迭代结构和角色定义，在处理运维和缺陷修复时使用看板的流程控制。例如，某河流水质监测系统的团队在新功能开发阶段采用Scrum方法，在系统上线后的运维阶段转为看板方法，既保证了开发过程的规范性，又提高了运维阶段的响应效率。

\subsection{2.2.3 极限编程（Extreme Programming,
XP）}\label{ux6781ux9650ux7f16ux7a0bextreme-programming-xp}

极限编程是一种注重技术实践的敏捷方法，强调高质量代码、持续反馈和团队协作。XP特别适合技术挑战高、质量要求严格的智慧水利平台核心模块开发。

\subsubsection{XP的核心价值观}\label{xpux7684ux6838ux5fc3ux4ef7ux503cux89c2}

XP基于五个核心价值观：

\begin{enumerate}
\def\labelenumi{\arabic{enumi}.}
\tightlist
\item
  \textbf{简单性}：

  \begin{itemize}
  \tightlist
  \item
    做最简单的能工作的设计
  \item
    避免过度设计和不必要的复杂性
  \item
    在智慧水利平台中，可能意味着先实现基本的水文数据分析功能，而非一开始就构建复杂的预测模型
  \end{itemize}
\item
  \textbf{沟通}：

  \begin{itemize}
  \tightlist
  \item
    强调团队成员之间的直接交流
  \item
    通过结对编程和站立会议促进沟通
  \item
    在智慧水利平台中，特别强调水利专家与开发人员之间的紧密沟通
  \end{itemize}
\item
  \textbf{反馈}：

  \begin{itemize}
  \tightlist
  \item
    尽早获取反馈并持续调整
  \item
    通过自动化测试、频繁集成和客户演示获取反馈
  \item
    在智慧水利平台中，可通过水利实测数据验证模型结果，不断调整算法
  \end{itemize}
\item
  \textbf{勇气}：

  \begin{itemize}
  \tightlist
  \item
    敢于重构代码和修改设计
  \item
    勇于承认和纠正错误
  \item
    在智慧水利平台中，可能需要勇于放弃不适合的技术路线，如发现某算法不适用时及时转向
  \end{itemize}
\item
  \textbf{尊重}：

  \begin{itemize}
  \tightlist
  \item
    团队成员互相尊重
  \item
    尊重客户和他们的专业知识
  \item
    在智慧水利平台中，尊重水利专家的领域知识和经验
  \end{itemize}
\end{enumerate}

\subsubsection{XP的核心实践}\label{xpux7684ux6838ux5fc3ux5b9eux8df5}

XP定义了12个核心实践，可分为四类：

\textbf{编程实践}： 1. \textbf{测试驱动开发（TDD）}： -
先编写测试，再编写代码 - 确保所有代码都有对应的测试 -
在智慧水利平台中，特别适用于水文计算模块等核心算法开发

\begin{enumerate}
\def\labelenumi{\arabic{enumi}.}
\setcounter{enumi}{1}
\tightlist
\item
  \textbf{结对编程}：

  \begin{itemize}
  \tightlist
  \item
    两名程序员共同工作在一台计算机上
  \item
    一人负责编码，一人负责审查和思考
  \item
    在智慧水利平台中，可采用''领域专家+开发者''的结对方式
  \end{itemize}
\item
  \textbf{简单设计}：

  \begin{itemize}
  \tightlist
  \item
    始终使用满足当前需求的最简设计
  \item
    避免不必要的复杂性和过早优化
  \item
    在智慧水利平台中，可先实现基础功能，再逐步优化性能
  \end{itemize}
\item
  \textbf{重构}：

  \begin{itemize}
  \tightlist
  \item
    在不改变行为的前提下改进代码结构
  \item
    持续提高代码可维护性
  \item
    在智慧水利平台中，随着需求变化和系统增长，定期重构数据处理管道等关键模块
  \end{itemize}
\end{enumerate}

\textbf{集成实践}： 5. \textbf{持续集成}： -
频繁集成和测试代码变更（通常每天多次） - 快速发现和修复集成问题 -
在智慧水利平台中，确保各个分析模块和数据处理组件能正确协同工作

\begin{enumerate}
\def\labelenumi{\arabic{enumi}.}
\setcounter{enumi}{5}
\tightlist
\item
  \textbf{共同代码所有权}：

  \begin{itemize}
  \tightlist
  \item
    任何团队成员都可以修改任何代码
  \item
    通过集体责任提高代码质量
  \item
    在智慧水利平台中，避免知识孤岛，确保核心功能有多人理解
  \end{itemize}
\end{enumerate}

\textbf{规划实践}： 7. \textbf{小型发布}： - 频繁发布小的功能增量 -
快速获取用户反馈 -
在智慧水利平台中，可以先发布简单的数据可视化功能，而非等待复杂分析模块全部完成

\begin{enumerate}
\def\labelenumi{\arabic{enumi}.}
\setcounter{enumi}{7}
\tightlist
\item
  \textbf{可持续节奏}：

  \begin{itemize}
  \tightlist
  \item
    保持稳定、可持续的工作速度
  \item
    避免长时间加班和倦怠
  \item
    在智慧水利平台中，特别要注意汛期等关键时段的工作安排
  \end{itemize}
\end{enumerate}

\textbf{团队实践}： 9. \textbf{客户在现场}： -
确保业务专家随时可用于解答问题 - 快速澄清需求和验证功能 -
在智慧水利平台中，可安排水利专家定期或者全职参与开发团队工作

\begin{enumerate}
\def\labelenumi{\arabic{enumi}.}
\setcounter{enumi}{9}
\tightlist
\item
  \textbf{代码规范}：

  \begin{itemize}
  \tightlist
  \item
    团队遵循统一的编码标准
  \item
    提高代码一致性和可读性
  \item
    在智慧水利平台中，特别关注水文模型参数和数据处理算法的标准化
  \end{itemize}
\item
  \textbf{集体所有权}：

  \begin{itemize}
  \tightlist
  \item
    所有团队成员共同对产品质量负责
  \item
    鼓励知识共享和协作
  \item
    在智慧水利平台中，避免单点依赖，确保关键知识在团队中共享
  \end{itemize}
\item
  \textbf{隐喻}：

  \begin{itemize}
  \tightlist
  \item
    使用共同的语言和比喻描述系统
  \item
    促进技术人员和业务人员之间的沟通
  \item
    在智慧水利平台中，建立水利概念与软件实现之间的映射关系
  \end{itemize}
\end{enumerate}

\subsubsection{XP在智慧水利平台中的应用}\label{xpux5728ux667aux6167ux6c34ux5229ux5e73ux53f0ux4e2dux7684ux5e94ux7528}

XP方法在智慧水利平台的以下场景中特别有价值：

\begin{enumerate}
\def\labelenumi{\arabic{enumi}.}
\tightlist
\item
  \textbf{水文模型和算法开发}：

  \begin{itemize}
  \tightlist
  \item
    通过TDD确保计算准确性和稳定性
  \item
    结对编程结合水利专家和程序员的优势
  \item
    案例：某流域洪水预报系统采用TDD开发关键水文计算模块，测试用例涵盖历史典型洪水场景，确保了模型在极端情况下的计算准确性
  \end{itemize}
\item
  \textbf{高可靠性系统开发}：

  \begin{itemize}
  \tightlist
  \item
    自动化测试和持续集成提高系统稳定性
  \item
    适用于防汛抗旱等关键业务系统
  \item
    案例：某大型水库调度系统采用结对编程和TDD开发决策支持模块，保证了系统在极端条件下的可靠运行
  \end{itemize}
\item
  \textbf{创新功能探索开发}：

  \begin{itemize}
  \tightlist
  \item
    小步快速迭代验证创新想法
  \item
    适用于智能预警、AR/VR水利应用等
  \item
    案例：某水利大数据分析平台在开发新型预警算法时采用XP方法，通过快速原型和持续反馈，在短时间内探索了多种算法组合
  \end{itemize}
\item
  \textbf{跨学科协作项目}：

  \begin{itemize}
  \tightlist
  \item
    通过结对工作和共同语言促进沟通
  \item
    适用于需要深度融合水利专业知识和IT技术的项目
  \item
    案例：某河流生态环境监测系统开发团队采用''生态专家+开发者''的结对模式，成功将复杂的生态评价模型转化为可操作的软件实现
  \end{itemize}
\end{enumerate}

\subsection{各敏捷方法的选择指南}\label{ux5404ux654fux6377ux65b9ux6cd5ux7684ux9009ux62e9ux6307ux5357}

在智慧水利平台开发中，应根据项目特点和团队情况选择合适的敏捷方法：

\begin{longtable}[]{@{}
  >{\raggedright\arraybackslash}p{(\columnwidth - 4\tabcolsep) * \real{0.2647}}
  >{\raggedright\arraybackslash}p{(\columnwidth - 4\tabcolsep) * \real{0.2647}}
  >{\raggedright\arraybackslash}p{(\columnwidth - 4\tabcolsep) * \real{0.4706}}@{}}
\toprule\noalign{}
\begin{minipage}[b]{\linewidth}\raggedright
项目特点
\end{minipage} & \begin{minipage}[b]{\linewidth}\raggedright
推荐方法
\end{minipage} & \begin{minipage}[b]{\linewidth}\raggedright
智慧水利平台实例
\end{minipage} \\
\midrule\noalign{}
\endhead
\bottomrule\noalign{}
\endlastfoot
需求变化频繁，需要规范流程 & Scrum &
智慧水利综合管理平台、防汛指挥系统 \\
持续性工作，需求稳定 & 看板 & 水利数据中心运维、监测设备管理系统 \\
技术挑战高，质量要求严格 & XP & 水文模型库、洪水预警决策支持系统 \\
大型复杂项目，多团队协作 & Scrum + 看板 &
流域综合管理平台（开发团队用Scrum，运维团队用看板） \\
创新探索性项目 & XP + Scrum &
智能水质预警系统（核心算法用XP，整体框架用Scrum） \\
\end{longtable}

无论选择哪种方法，都应遵循以下原则：

\begin{enumerate}
\def\labelenumi{\arabic{enumi}.}
\tightlist
\item
  \textbf{适应性调整}：根据团队反馈和项目进展调整方法实践
\item
  \textbf{关注价值交付}：选择能够最快交付价值的方法和实践
\item
  \textbf{团队参与决策}：让团队参与方法选择和调整过程
\item
  \textbf{渐进式采用}：从核心实践开始，逐步扩展和完善
\item
  \textbf{持续改进}：定期评估方法效果，不断优化流程
\end{enumerate}

\subsection{下一步学习}\label{ux4e0bux4e00ux6b65ux5b66ux4e60}

了解了主要敏捷方法论的理论和实践后，接下来我们将在\href{section03-02-03.md}{2.3
敏捷实践在智慧水利项目中的应用}中，详细探讨如何将这些方法具体应用于智慧水利平台的开发实践中，包括用户故事编写、迭代计划、团队组建和持续集成等关键环节。
